\section*{Appendix: Implementation and Reproducibility Details}
\addcontentsline{toc}{section}{Appendix}
\label{app:impl}

\paragraph{Hardware and Runtime.}
Qwen3-8B embeddings were generated on an NVIDIA RTX\,3090
(124{,}158 patients, $\approx$7\,h per embedding run).
Bio\_ClinicalBERT (110M params) ran on an NVIDIA RTX\,4090 at
$\approx$30 min per pass.
CoxPH models were fitted on CPU with \texttt{lifelines} v0.27;
evaluation used \texttt{scikit-survival} for TD-AUC and IBS.
All PCA projections and CoxPH coefficients are fitted on the training
split only; validation and test transforms use training-fit parameters.
\begin{sloppypar}Code and evaluation scripts:
{\small\url{https://github.com/Finley-Maple/Death\_Clock}}.\end{sloppypar}

\paragraph{Feature Dimensions.}
Table~\ref{tab:dims} summarises the exact dimensionalities at each stage of
the pipeline.
The 39-variable count refers to raw clinical variables before any categorical
encoding (demographics, 21 biomarkers, 11 binary comorbidity flags, sex,
smoking, alcohol status); this is the input width for S1 and the raw-side
of concatenation fusion.
For summation fusion, categorical variables are independently standardised
with the same 39-variable block; the PCA projection is matched to the same
39-dimensional width so that element-wise addition is dimension-consistent.

\begin{table}[H]
\centering
\caption{Feature and embedding dimensions by pipeline stage.}
\label{tab:dims}
\begin{tabular}{lll}
\toprule
Stage & Dimension & Description \\
\midrule
Raw feature block          & 39  & Demographics, biomarkers, comorbidities \\
PCA embedding (concat)     & 64  & Fit on training split only \\
CoxPH input (S1)           & 39  & Direct raw features \\
CoxPH input (concat)       & 103 & 39 raw $+$ 64 PCA \\
Summation fusion output    & 39  & PCA matched to raw feature width \\
\bottomrule
\end{tabular}
\end{table}
